\section{任务、问题与记号约定}
\label{sec:intro}

\subsection{本报告要回答的问题}

第一次讨论对本部分的原文要求是：

\begin{quote}\itshape
3. EMT（仅需了解，后续大概率不会用到）与多尺度建模（重要）。
EMT 模型最复杂，简单科普一下 EMT 即可。需要调研电网的多尺度建模：先了解
潮流/RMS/EMT 这 3 类模型的假设是什么、忽略了哪些物理过程、研究哪个时间尺度、
有哪些变量；然后调研潮流和 RMS 模型之间怎么连接、之间传什么变量等。
Kuramoto 模型为什么可以用来研究电网，用 Kuramoto 模型研究电网假设是什么、
并解释变量含义、Kuramoto 模型与 RMS 模型有什么关系。
\end{quote}

据此把任务拆成四个可以独立回答、也可以被逐条质疑的问题：

\begin{description}[leftmargin=2.2em]
\item[Q1（EMT 简介）] 什么是电磁暂态（EMT）模型？它保留了什么别的模型丢弃的物理过程？
为什么它不可替代，但在本项目中不作为主线？
\item[Q2（三类模型对照）] 潮流（PF/QSTS）、RMS、EMT 各自的数学形式、假设、
忽略的物理过程、时间尺度、状态变量是什么？三者之间是“不同精度的同一模型”，
还是“不同假设下的不同模型”？
\item[Q3（多尺度连接，重点）] 潮流与 RMS 之间如何连接、传什么变量？
为什么可以把完整模型降阶为相量模型，误差有多大？混合 EMT--RMS 仿真的接口
交换什么、代价是什么？
\item[Q4（Kuramoto，重点）] 为什么 Kuramoto 振子模型能研究电网？
它做了哪些假设、变量对应什么物理量？它与 RMS 模型是等价、近似还是降阶？
\end{description}

\subsection{记号与核心概念约定}
\label{sec:notation}

为了避免“同一个词在不同章节指不同东西”，先固定本报告的用法。
每个定义都给出\textbf{为什么需要它}，后文首次使用时不再重复解释。

\begin{description}[leftmargin=2.4em]
\item[模型] 指“一组方程 $+$ 一组假设 $+$ 一个适用范围”。因此说“EMT 模型”时，
同时意味着它的假设（保留全部 $L,C$ 动态、可含开关）和适用范围（$\mu$s--ms）。
本报告中“三类模型”始终指 PF/QSTS、RMS、EMT。
\item[状态变量] 指真正带时间导数、需要初值的量（记 $\bx$）；
与之相对的是由代数约束即时确定的量（记 $\bm{y}$，如潮流中的 $V,\theta$）。
区分二者是理解“PF 没有状态、RMS 有状态”的关键。
\item[平衡点（工作点）] 令所有时间导数为零后方程的解。潮流解就是
RMS 方程~\eqref{eq:rms-dae} 的平衡点；这是 \S\ref{sec:workpoint} 的第一层连接。
\item[相量] 把一个窄带正弦信号 $a(t)=A(t)\cos(\omega_s t+\varphi(t))$
的幅值与相位打包为复数 $\underline{A}=A e^{\jmath\varphi}$
（用 $A$ 而非有效值时也称“峰值相量”）。
相量成立的前提是 $A(t),\varphi(t)$ 相对载波（此处指 $50$ Hz 工频正弦）变化足够慢——这个“足够慢”
在 \S\ref{sec:quasi} 被量化为一个不等式。
\textbf{相量不是某一类模型专属}：PF 与 RMS 都用它，区别在于\textbf{相量是否随时间变化}——
PF 中相量取常数（因此是代数方程），RMS 中相量是随时间慢变的复包络（因此是微分方程）；
EMT 不用相量，直接求瞬时值。
\item[降阶] 在\textbf{关注的时间尺度上}保持所需精度的前提下，用更少的状态或更弱的方程替代原系统；
降阶的代价是丢掉该尺度之外的动态与非线性细节（因此“潮流丢掉全部动态”与“Kuramoto 保留相角动态”
都是合法的降阶，只是关注点不同）。本报告的三类模型与 Kuramoto 正处在这样一个简化关系之中（式~\eqref{eq:chain}）。
\item[时间尺度分离] 存在小参数 $0<\varepsilon\ll1$，使系统可写成
$\dot\bx=\bm f(\bx,\bz),\ \varepsilon\dot\bz=\bm g(\bx,\bz)$（$\bx$ 为慢变量、$\bz$ 为快变量）。
它使“忽略快变量”从经验做法变成有误差界可言的数学操作（\S\ref{sec:singular}）。
\item[多尺度建模] 指在同一研究目标下\textbf{同时或交替}使用不同尺度的模型，
并且\textbf{显式定义}它们之间交换什么变量、以什么步长、误差多大
（\S\ref{sec:q3} 的三层连接）。
\end{description}

\subsection{与项目书技术点的对应}

本部分对应项目书技术内容第一点“基于时空动力学的分布式资源高维非线性接入行为建模技术”，
以及技术难点第二点“多源外部不确定因素对电网运行影响难以准确刻画”，
并为技术内容第二点（有限元/伽辽金改进潮流计算）提供“哪些动态可以安全忽略”
的判据依据：只有先说清各模型忽略了什么，才能判断在什么条件下可以安全地降阶与加速。

\subsection{阅读路线}

若时间有限，只读三处即可覆盖全部结论：表~\ref{tab:retention}（三模型保留矩阵）、图~\ref{fig:chain}（模型之间的简化关系）、表~\ref{tab:assumptions}（每条结论的假设与失效条件）。
如需查阅具体步骤与公式，见正文的 6 个推导块（【推导块 1--6】，分别位于 \S\ref{sec:q3}、\S\ref{sec:q4}）；
符号与缩略语分别见附录~\ref{app:symbols}、附录~\ref{app:abbr}，设备与概念释义见附录~\ref{app:concepts}。
